\begin{exercise}[subtitle=Pauli-Prinzip, points=1, Punkt=1]
Erklären Sie das Pauli-Prinzip und seine Bedeutung für den Aufbau der Elektronenhülle.
\begin{solution}
Das Pauli-Prinzip besagt, dass zwei Elektronen in einem Atom nicht in allen vier Quantenzahlen übereinstimmen können. Folglich dürfen in einem Orbital maximal zwei Elektronen mit entgegengesetztem Spin sitzen. Dieses Prinzip erklärt die Elektronenkonfigurationen und die Struktur des Periodensystems.
\end{solution}
\end{exercise}

\begin{exercise}[subtitle=Bindungsarten, points=1, Punkt=1]
Beschreiben Sie den Unterschied zwischen einer Ionenbindung und einer kovalenten Bindung.
\begin{solution}
Bei einer Ionenbindung werden Elektronen von einem Atom zum anderen übertragen, wodurch positiv und negativ geladene Ionen entstehen, die durch elektrostatische Kräfte zusammengehalten werden. Eine kovalente Bindung entsteht hingegen durch das Teilen von Elektronenpaaren zwischen Atomen.
\end{solution}
\end{exercise}

\begin{exercise}[subtitle=Edelgase, points=1, Punkt=1]
Warum sind Edelgase (z.B. Helium, Neon) chemisch in der Regel inert?
\begin{solution}
Edelgase besitzen bereits eine vollständig gefüllte äußere Elektronenschale und erfüllen damit die Oktettregel. Dadurch haben sie kein Bestreben, Elektronen aufzunehmen oder abzugeben, was sie in den meisten Fällen reaktionsträge macht.
\end{solution}
\end{exercise}

\begin{exercise}[subtitle=Endotherm vs. exotherm, points=1, Punkt=1]
Erklären Sie den Unterschied zwischen endothermen und exothermen Reaktionen.
\begin{solution}
Exotherme Reaktionen setzen Energie (meist Wärme) an die Umgebung frei; ihre Produkte haben weniger Energie als die Reaktanten. Endotherme Reaktionen benötigen hingegen Energiezufuhr aus der Umgebung; die Produkte haben mehr Energie als die Ausgangsstoffe.
\end{solution}
\end{exercise}

\begin{exercise}[subtitle=pH-Wert, points=1, Punkt=1]
Definieren Sie den pH-Wert und geben Sie die pH-Skala mit ihren wichtigsten Bereichen an.
\begin{solution}
Der pH-Wert ist der negative dekadische Logarithmus der Wasserstoffionen-Konzentration: \( \mathrm{pH} = -\log_{10}[H_3O^+] \). Die Skala reicht typischerweise von 0 (stark sauer) über 7 (neutral) bis 14 (stark basisch).
\end{solution}
\end{exercise}

\begin{exercise}[subtitle=Avogadro-Gesetz, points=1, Punkt=1]
Was besagt das Avogadro-Gesetz im Kontext idealer Gase?
\begin{solution}
Das Avogadro-Gesetz besagt, dass bei konstantem Druck und konstanter Temperatur gleiche Volumina idealer Gase die gleiche Anzahl Teilchen (Mole) enthalten. Es führt zur Definition der Avogadro-Zahl.
\end{solution}
\end{exercise}

\begin{exercise}[subtitle=Isotope, points=1, Punkt=1]
Was sind Isotope eines Elements und wie unterscheiden sie sich voneinander?
\begin{solution}
Isotope sind Atome desselben Elements mit gleicher Protonenzahl (Ordnungszahl), aber unterschiedlicher Neutronenzahl. Sie besitzen daher unterschiedliche Massenzahlen und können sich in physikalischen Eigenschaften wie der Stabilität unterscheiden.
\end{solution}
\end{exercise}

\begin{exercise}[subtitle=Polarität, points=1, Punkt=1]
Warum ist das Wassermolekül polar?
\begin{solution}
Die Bindungen in Wasser sind polar, weil Sauerstoff eine höhere Elektronegativität als Wasserstoff besitzt. Zusätzlich ist die Molekülgeometrie gewinkelt (ca. 104{,}5°), sodass sich die Dipolmomente der O–H-Bindungen nicht gegenseitig aufheben. Dadurch entsteht ein Gesamt-Dipolmoment.
\end{solution}
\end{exercise}

\begin{exercise}[subtitle=Oxidation und Reduktion, points=1, Punkt=1]
Definieren Sie die Begriffe Oxidation und Reduktion im Sinne der Elektronenübertragung.
\begin{solution}
Oxidation ist die Abgabe von Elektronen, wodurch die Oxidationszahl steigt. Reduktion ist die Aufnahme von Elektronen, wodurch die Oxidationszahl sinkt. Beide Prozesse laufen gekoppelt ab.
\end{solution}
\end{exercise}

\begin{exercise}[subtitle=Le Chatelier, points=1, Punkt=1]
Was besagt das Prinzip von Le Chatelier (Prinzip des kleinsten Zwangs)?
\begin{solution}
Wenn auf ein im Gleichgewicht befindliches System ein Zwang ausgeübt wird (z. B. Änderung von Konzentration, Temperatur oder Druck), verschiebt sich das Gleichgewicht so, dass es dem Zwang entgegenwirkt, um einen neuen Gleichgewichtszustand zu erreichen.
\end{solution}
\end{exercise}

\begin{exercise}[subtitle=Brønsted-Lowry-Säure, points=1, Punkt=1]
Wie definiert die Brønsted-Lowry-Theorie Säuren und Basen?
\begin{solution}
Nach Brønsted-Lowry ist eine Säure ein Protonendonator, der ein Proton an eine Base abgeben kann. Eine Base ist ein Protonenakzeptor, der ein Proton von einer Säure aufnehmen kann.
\end{solution}
\end{exercise}

\begin{exercise}[subtitle=Aktivierungsenergie, points=1, Punkt=1]
Was versteht man unter Aktivierungsenergie und wie beeinflusst sie die Reaktionsgeschwindigkeit?
\begin{solution}
Die Aktivierungsenergie ist die Mindestenergie, die nötig ist, um eine Reaktion zu starten, indem die Reaktanten in einen Übergangszustand überführt werden. Je höher die Aktivierungsenergie, desto langsamer verläuft die Reaktion bei gegebener Temperatur.
\end{solution}
\end{exercise}

\begin{exercise}[subtitle=Atommasse vs. Molmasse, points=1, Punkt=1]
Erklären Sie den Unterschied zwischen der Atommasse und der Molmasse.
\begin{solution}
Die Atommasse bezieht sich auf die Masse eines einzelnen Atoms, ausgedrückt in atomaren Masseneinheiten (u). Die Molmasse ist die Masse eines Mols einer Substanz, angegeben in Gramm pro Mol (g/mol), und entspricht numerisch der relativen Molekülmasse.
\end{solution}
\end{exercise}

\begin{exercise}[subtitle=Hybridisierung, points=1, Punkt=1]
Was ist Hybridisierung in der Chemie, und warum ist sie für die Molekülgeometrie wichtig?
\begin{solution}
Hybridisierung ist das Mischen von Atomorbitalen (z. B. s- und p-Orbitalen) zu gleichwertigen Hybridorbitalen. Diese bestimmen die räumliche Anordnung der Bindungen und erklären z. B. die tetraedrische Geometrie von Methan (sp\(^3\)-Hybridisierung).
\end{solution}
\end{exercise}

\begin{exercise}[subtitle=Zustandsgleichungen, points=1, Punkt=1]
Nennen Sie die allgemeine Zustandsgleichung für ideale Gase und erklären Sie die Bedeutung der Symbole.
\begin{solution}
Die allgemeine Zustandsgleichung lautet \( pV = nRT \). Dabei ist \( p \) der Druck, \( V \) das Volumen, \( n \) die Stoffmenge, \( R \) die allgemeine Gaskonstante und \( T \) die absolute Temperatur. Sie beschreibt den Zusammenhang dieser Größen für ein ideales Gas.
\end{solution}
\end{exercise}
